\documentclass[12pt, a4paper]{article}

\usepackage[top=2.5cm, bottom=2.5cm, left=3cm, right=3cm]{geometry}
\usepackage{mdframed}
\usepackage[utf8]{inputenc}
\usepackage[T1]{fontenc}
\usepackage[spanish]{babel}
\usepackage{amsmath}
\usepackage{float}
\usepackage[margin=10pt, font=small, labelfont=bf]{caption}
\usepackage{booktabs}
\usepackage{graphicx}
\usepackage{xcolor}
\usepackage{hyperref}
\usepackage{natbib}
\usepackage{longtable}
\usepackage{comment}
\usepackage{setspace}
\usepackage{titlesec}
\usepackage{fancyhdr}

\setlength{\textfloatsep}{18pt plus 2pt minus 4pt}
\setlength{\floatsep}{12pt plus 2pt minus 2pt}
\setlength{\intextsep}{14pt plus 2pt minus 2pt}
\setstretch{1.15}
\setlength{\parindent}{0pt}
\setlength{\parskip}{8pt}

\titleformat{\section}
  {\large\bfseries}{\thesection.}{0.5em}{}
  [\vspace{2pt}\hrule\vspace{4pt}]
\titleformat{\subsection}
  {\normalsize\bfseries}{\thesubsection.}{0.5em}{}
\titlespacing{\section}{0pt}{18pt}{8pt}
\titlespacing{\subsection}{0pt}{14pt}{6pt}

\hypersetup{colorlinks=true, linkcolor=blue!60!black}

\providecommand{\reportnumber}{}
\providecommand{\reporttitle}{}
\providecommand{\reportsubtitle}{}

\pagestyle{fancy}
\fancyhf{}
\fancyhead[L]{\footnotesize Centro Javeriano de Competitividad}
\fancyhead[R]{\footnotesize \reportnumber}
\cfoot{\thepage}
\setlength{\headheight}{14pt}
\renewcommand{\headrulewidth}{0.4pt}

\geometry{margin=1in}
\graphicspath{{Paper/figures/}{figures/}}
\newcommand{\tablepath}{Paper/tables/}
\newcommand{\figurelang}{es}
\newcommand{\reportseries}{Informes sobre la Productividad en Colombia}
\newcommand{\reportseriesdesc}{Evidencia para entender y fomentar la productividad}
\newcommand{\reportauthorone}{Oliver Pardo}
\newcommand{\reportauthoroneaffil}{Profesor asociado del Departamento de Administración, Pontificia Universidad Javeriana; director del Centro Javeriano de Competitividad}
\newcommand{\reportauthortwo}{Jorge Orozco}
\newcommand{\reportauthortwoaffil}{Consultor de investigación del Centro Javeriano de Competitividad, Pontificia Universidad Javeriana}
\renewcommand{\contentsname}{Tabla de contenido}
\newif\ifpolicyreport
\policyreporttrue


\renewcommand{\reportnumber}{Informe 03}
\renewcommand{\reporttitle}{La Geografía de la Remuneración Laboral en Colombia}
\renewcommand{\reportsubtitle}{Diferencias departamentales, 2010--2024pr}

\title{\reporttitle\\\reportsubtitle}
\author{\reportauthorone \and \reportauthortwo}
\date{\today}

\begin{document}

\begin{titlepage}
\thispagestyle{empty}
\vspace*{0.6cm}

\begin{center}

\includegraphics[width=0.58\textwidth]{Paper/figures/logo_cjc_palatino_horizontal.png}
\end{center}

\vspace{1.5cm}

\begin{center}
{\Large\bfseries \reportseries\par}
\vspace{0.15cm}
{\large \reportseriesdesc\par}

\vspace{1.4cm}

{\large\bfseries \reportnumber\par}
\vspace{0.45cm}
{\Huge\bfseries \reporttitle\par}
\vspace{0.35cm}
{\Large \reportsubtitle\par}
\vspace{0.65cm}
{\large\bfseries \reportauthorone\par}
{\normalsize \reportauthoroneaffil\par}
\vspace{0.25cm}
{\large\bfseries \reportauthortwo\par}
{\normalsize \reportauthortwoaffil\par}
\end{center}

\vfill

\begin{center}
{\large \today\par}
\end{center}

\end{titlepage}


\begin{mdframed}[
  linewidth=1pt,
  linecolor=gray!60,
  backgroundcolor=gray!6,
  innertopmargin=12pt,
  innerbottommargin=12pt,
  innerleftmargin=14pt,
  innerrightmargin=14pt
]
\section*{Resumen ejecutivo}
\addcontentsline{toc}{section}{Resumen ejecutivo}

\textbf{Este informe analiza la remuneración laboral por trabajador en los departamentos de Colombia.} El ejercicio usa la GEIH para construir una remuneración mensual promedio por ocupado, expresada en pesos constantes de 2025, para los 24 departamentos comparables durante todo el periodo.

\textbf{Bogotá D.C. ocupa una posición claramente dominante en niveles de remuneración.} En 2024pr registra una remuneración mensual por trabajador cercana a 2,95 millones de pesos de 2025. Le siguen Antioquia, Caldas, Meta, Risaralda, Santander y Valle del Cauca. En el extremo inferior aparecen La Guajira, Sucre, Córdoba, Nariño y Cauca.

\textbf{El crecimiento de la remuneración también fue desigual.} Risaralda, Cundinamarca, Caldas y Bogotá D.C. registraron los mayores crecimientos anualizados entre 2010 y 2024pr. En contraste, La Guajira, Cesar, Magdalena y Sucre estuvieron entre los desempeños más débiles.

\textbf{La remuneración laboral tiene una dimensión territorial marcada.} Buena parte de los departamentos rezagados se ubica en la periferia norte, pacífica y sur del país, mientras que varios de los aceleradores se concentran en el centro andino. Esta distribución sugiere que la remuneración depende no solo de características individuales de los trabajadores, sino también de conectividad, acceso a mercados, infraestructura, densidad empresarial, composición productiva y capacidad institucional.

\textbf{La geografía de la remuneración prepara el terreno para estudiar su relación con la productividad.} Los departamentos con mayor remuneración no siempre coinciden con los más productivos, y los departamentos donde más crece la remuneración no necesariamente son aquellos donde más crece la productividad.
\end{mdframed}

\newpage
\tableofcontents
\newpage

\section{Introducción}

\textbf{La remuneración laboral también tiene una geografía.} Las diferencias de ingreso entre trabajadores no dependen únicamente de su educación, ocupación o sector económico. También dependen del territorio donde trabajan, de la estructura productiva local, de la informalidad, de la conectividad con mercados y de la densidad empresarial de cada departamento.

\textbf{Este informe estudia la evolución de la remuneración por trabajador en los departamentos de Colombia entre 2010 y 2024pr.} El ejercicio se concentra en los 24 departamentos que pueden seguirse de manera comparable en la GEIH durante todo el periodo. La pregunta central es sencilla: ¿en qué departamentos se remunera más el trabajo y en cuáles ha crecido más la remuneración real?

\textbf{La lectura territorial de la remuneración es importante para la agenda de productividad.} Un aumento de la productividad solo mejora de manera amplia el bienestar si se traduce, al menos parcialmente, en mejores ingresos laborales. Antes de estudiar esa relación de manera directa, conviene entender primero la geografía propia de la remuneración.

\textbf{El informe está organizado en cinco secciones, incluyendo esta introducción.} La sección 2 presenta la metodología. La sección 3 compara los niveles de remuneración departamental. La sección 4 analiza los crecimientos y clasifica los departamentos según nivel y dinamismo de la remuneración. La sección 5 presenta conclusiones.

\section{Metodología}

\textbf{La medición utiliza los microdatos armonizados de la GEIH.} Para cada departamento y año se calcula la remuneración laboral mensual total a partir del ingreso laboral por hora, las horas semanales trabajadas y el factor de expansión de la encuesta. La remuneración se expresa en pesos constantes de 2025.

\textbf{La remuneración por trabajador se construye de forma comparable con el PIB por trabajador.} La remuneración laboral mensual total se calcula entre los ocupados con ingreso horario positivo y horas válidas. Luego, ese total se divide por el número total de ocupados del departamento. De esta manera, el indicador mide la remuneración mensual promedio por ocupado del departamento, manteniendo un denominador comparable con el usado en los ejercicios de productividad.

\textbf{El informe se concentra en 24 departamentos comparables durante todo el periodo.} Para evitar imputaciones o cambios artificiales en la cobertura, se excluyen San Andrés y Providencia y los departamentos de la Amazonía y la Orinoquía que no hacen parte del panel comparable de la GEIH desde 2010. Bogotá D.C. se trata como departamento porque así aparece en las cuentas nacionales departamentales.

\textbf{El periodo de análisis es 2010--2024pr y se excluye 2020.} El año 2020 se excluye para mantener consistencia con el tratamiento del Informe 01 y con el informe departamental de productividad. Las tasas de crecimiento reportadas son tasas anualizadas entre 2010 y 2024pr.

\section{Niveles de remuneración por departamento}

\textbf{El ranking de niveles de remuneración confirma la centralidad de Bogotá D.C., pero también muestra diferencias importantes entre departamentos.} El Cuadro \ref{tab:pib_geih_productividad_departamento_remuneracion_niveles} ordena los departamentos por remuneración por trabajador en 2024pr. Bogotá D.C. ocupa el primer lugar, seguida por Antioquia, Caldas, Meta, Risaralda, Santander y Valle del Cauca.

\begin{table}[H]
\centering
\caption{Ranking departamental de niveles de remuneraci\'on por trabajador, 2024}
\label{tab:pib_geih_productividad_departamento_remuneracion_niveles}
\scriptsize
\begin{tabular}{lrrrr}
\toprule
Departamento & Rem./trab. 2024 & Puesto & PIB/trab. 2024pr & Puesto \\
\midrule
Bogotá D.C. & 2,95 & 1 & 67,4 & 1 \\
Antioquia & 2,12 & 2 & 47,7 & 5 \\
Caldas & 1,85 & 3 & 36,6 & 12 \\
Meta & 1,83 & 4 & 63,1 & 2 \\
Risaralda & 1,82 & 5 & 38,7 & 10 \\
Santander & 1,80 & 6 & 62,0 & 3 \\
Valle del Cauca & 1,79 & 7 & 54,5 & 4 \\
Cundinamarca & 1,79 & 8 & 40,4 & 8 \\
Quindío & 1,71 & 9 & 33,8 & 13 \\
Atlántico & 1,58 & 10 & 37,3 & 11 \\
Tolima & 1,50 & 11 & 40,0 & 9 \\
Norte de Santander & 1,49 & 12 & 24,3 & 20 \\
Boyacá & 1,49 & 13 & 47,5 & 6 \\
Caquetá & 1,47 & 14 & 24,9 & 18 \\
Huila & 1,36 & 15 & 31,6 & 15 \\
Bolívar & 1,33 & 16 & 45,1 & 7 \\
Cesar & 1,26 & 17 & 32,6 & 14 \\
Chocó & 1,16 & 18 & 28,6 & 16 \\
Magdalena & 1,15 & 19 & 23,9 & 21 \\
Cauca & 1,06 & 20 & 26,6 & 17 \\
Nariño & 1,06 & 21 & 17,8 & 24 \\
Córdoba & 1,04 & 22 & 22,0 & 22 \\
Sucre & 1,01 & 23 & 22,0 & 23 \\
La Guajira & 0,89 & 24 & 24,8 & 19 \\
\bottomrule
\end{tabular}
\caption*{\footnotesize Nota: remuneraci\'on por trabajador en millones de pesos constantes de 2025 al mes; PIB por trabajador en millones de pesos constantes de 2015. La tabla est\'a ordenada por el nivel de remuneraci\'on por trabajador. Fuente: c\'alculos propios con DANE y GEIH.}
\end{table}

\textbf{La brecha entre los departamentos de mayor y menor remuneración es amplia.} En 2024pr, Bogotá D.C. registra una remuneración mensual por trabajador cercana a 2,95 millones de pesos de 2025. En el extremo inferior aparecen La Guajira, Sucre, Córdoba, Nariño y Cauca, con niveles cercanos o inferiores a 1,1 millones. La diferencia es suficientemente grande como para mostrar que el mercado laboral colombiano no puede entenderse solo desde el promedio nacional.

\textbf{Los departamentos con mayor remuneración no son necesariamente los mismos que tienen mayor productividad.} Antioquia y Caldas ocupan posiciones altas en remuneración, aunque no están en los primeros lugares de PIB por trabajador. En sentido contrario, Meta y Santander tienen niveles altos de PIB por trabajador, pero su posición en remuneración es menos favorable. Esta diferencia anticipa que la remuneración laboral depende de más factores que la producción promedio por ocupado.

\section{Crecimiento y geografía de la remuneración}

\textbf{El crecimiento de la remuneración por trabajador también fue desigual entre departamentos.} El Cuadro \ref{tab:pib_geih_productividad_departamento_remuneracion_crecimientos} muestra que Risaralda, Cundinamarca, Caldas y Bogotá D.C. registraron los mayores crecimientos anualizados de remuneración por trabajador entre 2010 y 2024pr. En el extremo contrario se ubican La Guajira, Cesar, Magdalena y Sucre.

\begin{table}[H]
\centering
\caption{Crecimiento departamental de remuneraci\'on y PIB por trabajador, 2010--2024pr}
\label{tab:pib_geih_productividad_departamento_remuneracion_crecimientos}
\scriptsize
\begin{tabular}{lrrrrr}
\toprule
Departamento & Crec. rem./trab. & Puesto & Crec. PIB/trab. & Puesto & Dif. puestos \\
\midrule
Risaralda & 2,83\% & 1 & 3,14\% & 3 & -2 \\
Cundinamarca & 2,39\% & 2 & -0,32\% & 23 & -21 \\
Caldas & 2,30\% & 3 & 2,14\% & 6 & -3 \\
Bogotá D.C. & 1,90\% & 4 & 1,98\% & 8 & -4 \\
Boyacá & 1,66\% & 5 & 1,16\% & 15 & -10 \\
Antioquia & 1,65\% & 6 & 1,27\% & 13 & -7 \\
Tolima & 1,55\% & 7 & 2,04\% & 7 & +0 \\
Quindío & 1,50\% & 8 & 3,78\% & 1 & +7 \\
Valle del Cauca & 1,47\% & 9 & 2,15\% & 5 & +4 \\
Meta & 1,43\% & 10 & 0,05\% & 20 & -10 \\
Bolívar & 1,20\% & 11 & 1,69\% & 10 & +1 \\
Cauca & 1,20\% & 12 & 1,44\% & 12 & +0 \\
Atlántico & 1,09\% & 13 & 0,46\% & 18 & -5 \\
Chocó & 1,09\% & 14 & 2,95\% & 4 & +10 \\
Norte de Santander & 1,06\% & 15 & 1,23\% & 14 & +1 \\
Caquetá & 0,97\% & 16 & 3,71\% & 2 & +14 \\
Huila & 0,95\% & 17 & -0,13\% & 22 & -5 \\
Córdoba & 0,79\% & 18 & 0,90\% & 17 & +1 \\
Nariño & 0,63\% & 19 & 0,97\% & 16 & +3 \\
Santander & 0,62\% & 20 & 1,71\% & 9 & +11 \\
Sucre & 0,35\% & 21 & 1,67\% & 11 & +10 \\
Magdalena & 0,24\% & 22 & 0,16\% & 19 & +3 \\
Cesar & 0,16\% & 23 & -0,03\% & 21 & +2 \\
La Guajira & -1,84\% & 24 & -0,80\% & 24 & +0 \\
\bottomrule
\end{tabular}
\caption*{\footnotesize Nota: tasas de crecimiento anualizadas. La diferencia de puestos corresponde al puesto en crecimiento de la remuneraci\'on por trabajador menos el puesto en crecimiento del PIB por trabajador; un valor positivo indica que el departamento ocupa una posici\'on m\'as baja en crecimiento de remuneraci\'on que en crecimiento de productividad. Fuente: c\'alculos propios con DANE y GEIH.}
\end{table}

\textbf{La comparación entre niveles y crecimientos permite distinguir cuatro situaciones.} La Figura \ref{fig:geih_remuneracion_departamento_cuadrantes} cruza el nivel de remuneración por trabajador en 2024pr con su crecimiento anualizado entre 2010 y 2024pr. La línea vertical muestra el crecimiento agregado y la línea horizontal muestra el nivel agregado de remuneración por trabajador.

\begin{figure}[H]
  \centering
  \includegraphics[width=\textwidth]{Paper/figures/fig_geih_remuneracion_departamento_cuadrantes.png}
  \caption{Nivel y crecimiento de la remuneración por trabajador por departamento, 2010--2024pr}
  \label{fig:geih_remuneracion_departamento_cuadrantes}
  \caption*{\footnotesize Nota: la línea vertical muestra el crecimiento agregado de la remuneración por trabajador y la línea horizontal muestra el nivel agregado de remuneración por trabajador en 2024pr. El tamaño de la burbuja es proporcional al número de ocupados del departamento. Fuente: cálculos propios con GEIH.}
\end{figure}

\textbf{Bogotá D.C. y Caldas combinan niveles altos de remuneración con crecimientos superiores al agregado.} Antioquia, Santander y Valle del Cauca mantienen remuneraciones altas, pero con crecimientos por debajo del agregado. En cambio, Cundinamarca, Meta y Risaralda aparecen como aceleradores de remuneración: no están por encima del nivel agregado en 2024pr, pero registraron aumentos superiores al promedio durante el periodo.

\begin{figure}[H]
  \centering
  \includegraphics[width=0.82\textwidth]{Paper/figures/fig_geih_remuneracion_departamento_mapa_cuadrantes.png}
  \caption{Mapa de cuadrantes de remuneración por trabajador, 2010--2024pr}
  \label{fig:geih_remuneracion_departamento_mapa_cuadrantes}
  \caption*{\footnotesize Nota: las categorías son las mismas de la Figura \ref{fig:geih_remuneracion_departamento_cuadrantes}. Los departamentos no comparables aparecen en gris. Fuente: cálculos propios con GEIH; geometría departamental de GADM.}
\end{figure}

\textbf{El mapa de cuadrantes muestra que la dinámica de la remuneración laboral tiene una dimensión territorial marcada.} Buena parte de los departamentos rezagados se ubica en la periferia norte, pacífica y sur del país, mientras que varios de los aceleradores se concentran en el centro andino. Esta distribución sugiere que las diferencias de remuneración no responden únicamente a características individuales de los trabajadores, sino también a condiciones territoriales más amplias: conectividad, acceso a mercados, infraestructura, densidad empresarial, composición productiva y capacidad institucional.

\section{Conclusiones}

\textbf{La remuneración laboral en Colombia presenta brechas territoriales profundas.} Incluso al limitar el análisis a los 24 departamentos comparables de la GEIH, la diferencia entre los departamentos de mayor y menor remuneración es amplia. Bogotá D.C. ocupa una posición claramente dominante, mientras que varios departamentos de la periferia tienen remuneraciones muy inferiores.

\textbf{El crecimiento de la remuneración no ha sido homogéneo.} Risaralda, Cundinamarca, Caldas y Bogotá D.C. registraron aumentos relativamente altos de remuneración por trabajador. La Guajira, Cesar, Magdalena y Sucre, en cambio, se ubicaron entre los desempeños más débiles. Esto implica que las brechas territoriales de remuneración no solo dependen de los niveles iniciales, sino también de trayectorias recientes muy distintas.

\textbf{La geografía de la remuneración debe ser parte de la agenda económica nacional.} La concentración de rezagos en la periferia norte, pacífica y sur del país sugiere que las políticas laborales y productivas necesitan una lectura territorial explícita. Mejorar la remuneración laboral no depende únicamente de reglas salariales nacionales; también requiere mejorar las condiciones que permiten a los territorios conectar trabajadores y empresas con actividades de mayor valor agregado.

\textbf{Este informe prepara el terreno para estudiar la relación entre productividad y remuneración.} Los departamentos con mayor remuneración no siempre coinciden con los departamentos más productivos, y los departamentos con mayor crecimiento de la remuneración no necesariamente son los de mayor crecimiento de productividad. Esa relación se analiza de manera directa en el siguiente informe de la serie.


\end{document}
